\section{Experimental Protocol}


This section describes the experimental procedure used to evaluate different
prompting strategies for algorithmic problem solving.



\subsection{Benchmark Preparation}


The experiment uses a curated benchmark dataset containing LeetCode-style
algorithm problems.


Each problem includes:

\begin{itemize}

\item Problem description.

\item Difficulty level.

\item Algorithm category.

\item Expected optimal approach.

\item Test cases.

\end{itemize}



The dataset covers multiple algorithm categories including arrays, hash tables,
linked lists, trees, graphs, and dynamic programming.



\subsection{LLM Solution Generation}


For each benchmark problem, the selected Large Language Model generates
solutions under three different prompt conditions:


\begin{enumerate}

\item Baseline Prompt.

\item Optimized Prompt.

\item AI-Assisted Workflow Prompt.

\end{enumerate}


The same model configuration and generation parameters are used for all
experiments to ensure fair comparison.



\subsection{Automated Evaluation Pipeline}


The evaluation pipeline follows the following procedure:


\begin{center}

Problem Input

$\rightarrow$

Prompt Generation

$\rightarrow$

LLM Response

$\rightarrow$

Code Extraction

$\rightarrow$

Program Execution

$\rightarrow$

Metric Calculation

\end{center}



Generated programs are automatically executed against benchmark test cases.



\subsection{Experiment Repetition}


To reduce randomness caused by LLM generation, experiments may be repeated
multiple times for each prompt strategy.


The final result is calculated using the average performance across runs.



\subsection{Result Collection}


For each experiment instance, the system records:


\begin{itemize}

\item Problem identifier.

\item Prompt strategy.

\item Generated solution.

\item Correctness score.

\item Runtime performance.

\item Complexity evaluation.

\end{itemize}



The collected results are stored in structured CSV and JSON formats for further
analysis and visualization.



\subsection{Reproducibility}


To enable reproducibility, this project provides:


\begin{itemize}

\item Dataset format.

\item Prompt templates.

\item Experiment scripts.

\item Evaluation framework.

\item Result processing tools.

\end{itemize}


The complete implementation will be publicly available through GitHub.
